
%\needspace{4\baselineskip}
\color{uniblue}{\subsection{Автоматическое повторное включение (АПВ)}}
\color{black}

\begin{enumerate}

\item
Около 85\% повреждений изоляции воздушных линий связано с дуговыми короткими замыканиями, которые по своей природе имеют неустойчивый характер и исчезают при отключении линии. Это означает, что линия может быть снова включена автоматикой повторного включения. В устройстве автоматическое повторное включение представлено функцией «АПВ».

\item
Логика АПВ реализует следующие возможности:
\begin{itemize}
    \item два цикла автоматического повторного включения после аварийного отключения выключателя (программный переключатель SGF4);
    \item формирование блокировок АПВ;
    \item формирование сигнала готовности АПВ;
    \item контроль напряжения шин и присоединения для пуска АПВ;
    \item пуск алгоритма улавливания синхронизма по факту срабатывания АПВ.
\end{itemize}

\item
Готовность циклов АПВ будет сформирована по истечению выдержки времени таймера Т1 при оперативном включении выключателя (наличие активного сигнала «В: Включен»), отсутствие воздействий на таймер Т6, а также при отсутствии сигналов блокировки включения выключателя от логики КСВ и запрета АПВ.

\item
Пусковые условия АПВ первого цикла формируются при выполнении следующих условий:
\begin{itemize}
    \item аварийное отключение выключателя; 
    \item готовность АПВ для данного цикла; 
    \item возможность проверки синхронизма (программный переключатель SGF3);
    \item контроль напряжения (программный переключатель SGF2).
\end{itemize}

\item
Команда на включение выключателя от первого цикла создается после завершения выдержки времени таймера Т2 при условии отсутствия блокировки включения выключателя.

\item
Если затянулась команда включения выключателя, то по истечении времени таймера Т3 сформируется выходной сигнал задержки включения выключателя «АПВ: Задержка включения» и сбросится готовность циклов АПВ.

\item
После завершения первого цикла АПВ, происходит сброс готовности этого цикла. Если выключатель снова отключился, и условия логики соблюдены, активируется второй цикл АПВ.

\item
По окончанию работы второго цикла АПВ происходит сброс готовности данного цикла и запускается таймер Т4, по истечении времени которого, восстанавливается готовность первого и второго цикла АПВ.

\item
Пусковые условия АПВ второго цикла формируются при наличии аварийного отключения выключателя, наличие готовности АПВ данного цикла с отсутствием готовности АПВ первого цикла и возможностью контроля синхронизма уставкой SGF3 с контролем напряжений уставкой SGF2. Команда включения выключателя от второго цикла формируется по истечению выдержки времени таймера Т5 с отсутствием блокировки включения выключателя.

\item
Сброс готовности первого цикла формируется при наличии одного из следующих условий:
\begin{itemize}
    \item функция АПВ выведена из работы;
    \item отсутствует сигнал РФК из логики КСВ или активны сигналы запрета АПВ как внешний, так и из внутренней логики;
    \item наличие первого цикла АПВ;
    \item сформирован сигнал задержки включения выключателя.
\end{itemize}

\item
Сброс готовности второго цикла формируется при наличии одного из следующих условий:
\begin{itemize}
    \item уставкой SGF4 второй цикл АПВ выведен из работы;
    \item функция «АПВ» выведена из работы;
    \item отсутствует сигнал РФК из логики КСВ или активны сигналы запрета АПВ как внешний, так и из внутренней логики;
    \item введен ключ блокировки 2 цикла «АПВ Блок. 2 цикла»;
    \item наличие срабатывания защит от ОЗЗ (ГСОЗЗ, НЗОЗЗ), при введенном положении уставки SGF5 и наличии второго цикла АПВ;
    \item по истечению выдержки времени таймера Т3 сформировался сигнал задержки включения выключателя;
    \item наличие второго цикла АПВ.
\end{itemize}

\item
Ввод функции в работу на этапе параметрирования устройства осуществляется программным переключателем SGF1 «Ввод\_функции». Информация о введенной функции отображается выходным сигналом «АПВ: Ввод».

\item
Функция может быть выведена из работы оперативно путем активации сигнала «ОВ АПВ», что характеризуется активным состоянием сигнала «АПВ: Оперативный вывод» на выходе функции.

\item
Алгоритм работы функции выполнен в соответствии с рисунком \ref{apv:img1}. В таблице \ref{apv:tbl1} приведены параметры, необходимые для настройки функции.

\vspace{1mm}
\begin{figure}[H]
\centering
\includegraphics[width=1\textwidth,height=1\textheight,keepaspectratio]{\fbpath/121. АПВ/_latex/img/apv.pdf}
\captionof{figure}{Функциональная схема логики АПВ}\label{apv:img1}
\end{figure}

\end{enumerate}
